\chapter{Putting it all together: one atmosphere, one timestep}\label{putting-it-together}

The preceding chapters describe individual equations and numerical choices.
A simulation becomes coherent only when those choices describe the same air,
water, geometry and time interval. This chapter follows a complete update and
shows how to reason about its coupled behavior. The main example is the
Cartesian C-grid split-explicit path with bulk microphysics. Other admitted
integrators and geometries have their own stage owners; the example is not a
claim that one ordering can be substituted for every supported configuration.

\section{Begin with a complete mathematical state}

A useful abstract state is
\[
 X=(\rho_d,\theta,\boldsymbol u,\boldsymbol q,\boldsymbol N,
       \boldsymbol Z,\boldsymbol s,\mathcal G,\mathcal H).
\]
Here $\boldsymbol q$ contains the selected water species, $\boldsymbol N$
contains the applicable number concentrations, and $\boldsymbol Z$ denotes
additional prognostic moments only where a scheme actually owns them.
$\boldsymbol s$ includes scheme-private physical state such as held column
sources or turbulence variables. $\mathcal G$ is geometry and domain placement;
$\mathcal H$ includes histories needed by boundaries, staged integration and
continuation. Some of these components are absent for a simpler configuration.
Pressure is additionally stored and used, but its committed relationship to
dry density and potential temperature is constrained by the production EOS.

The state is larger than the arrays shown in a plot. Two runs with identical
visible temperature and water at one instant can evolve differently if their
particle number, held heating interval, radiative boundary history or nested
placement differs. A restart specification must therefore identify the full
continuation state. A snapshot of plotted fields is insufficient evidence that
the next tendency will match.

Geometry is included because changing cell volume changes the interpretation
of density and mixing ratio. Time history is included because a source held
between radiation calls is not equivalent to a source freshly evaluated from
the current cloud. A complete state definition is the first defense against
accidentally comparing different mathematical problems while believing that
only a backend, thread count or file format changed.

\section{Initialization must create a mutually compatible state}

A sounding supplies a vertical environment, but the solver must place it on
its own scalar and velocity locations. The procedure needs physical heights,
a pressure/base-state construction, water conventions and an equation of state.
An interpolation that is monotone in a sounding coordinate does not
necessarily preserve hydrostatic balance on the model interfaces. Similarly,
using moisture in a thermodynamic conversion does not justify adding it a
second time to a buoyancy correction.

For a simple resting dry column, hydrostatic balance is
$dp_0/dz=-\rho_0g$. A discrete balance might be expressed at a face as
\[
 \frac{p_{0,k+1}-p_{0,k}}{z_{k+1}-z_k}
 =-g\rho_{0,k+1/2},
\]
provided this is the same averaging and distance used by the force operator.
That last condition is essential. Integrating a continuous hydrostatic profile
very accurately and sampling it with an incompatible force stencil can still
launch adjustment waves. Nimbus's initialization and hydrostatic tests examine
the implemented pairing; they are more informative than assuming that the word
``hydrostatic'' in a configuration guarantees a discrete resting solution.

A thermal perturbation is also a declared intervention in the initial state.
Holding pressure fixed while increasing theta changes diagnosed temperature
and density differently from holding density fixed. A bubble specification
must be interpreted through the actual initialization path, not through an
informal statement that it ``adds heat.'' The
\href{../src/init}{initialization implementation},
\href{../tests/init}{initialization tests} and the public configuration compiler
make those choices concrete. Invalid dimensions, missing coefficient tables,
incompatible boundaries and unavailable combinations must be rejected before
advancing a partially constructed state.

\section{Choose the interval before assigning its sources}

The outer runtime chooses an allowed large time step from configuration, current
state and the relevant stability constraints. The acoustic substep count then
belongs to the chosen split method. Output times, source cadences and a requested
end time may further bound an interval. A process that receives a smaller
interval must use that interval consistently for its tendency integration,
precipitation and diagnostic receipts.

For a source evaluated every $\tau$ seconds, a held tendency can be written
schematically as
\[
 S(t)=S(X(t_m)),\qquad t_m\le t<t_m+\tau.
\]
This is an approximation to the continuously reevaluated source. Its error
depends on how rapidly the source and its inputs vary during the hold, not
merely on whether $\tau$ divides a nominal large step. For example, rapidly
changing condensate can change radiative heating while the held value remains
unchanged. A continuation must preserve both the held value and its timing.
Reevaluating it immediately after restart changes the experiment.

Nested levels add a synchronization problem: a child's shorter steps cover a
parent interval, but a parent source or boundary history is not automatically
valid at each child stage. The execution plan must specify interpolation,
refresh and feedback times. The public admission system exists partly because
individual schemes being selectable does not establish those combined owners.

\section{Transport dry air and its carried quantities coherently}

For the representative split path, slow tendencies are prepared from the
appropriate stage state. The integrator applies the slow increment, seeds the
acoustic radiation history, and advances the fast pressure/density and momentum
sequence through its acoustic substeps. Closure of the fields read by the next
substep is part of that sequence. Reusing a host or device buffer is correct
only if it holds the state required at this particular stage.

The continuity calculation produces a carrier mass change. Retained face fluxes
allow the subsequent scalar update to use the same mass that moved through each
face. For a species $x$, the coupled relation is
\[
 \begin{aligned}
 M_C^*&=M_C^n-\sum_f B_{Cf}\mathcal F_f,\\
 W_{x,C}^*&=M_C^n q_{x,C}^n-\sum_fB_{Cf}\mathcal F_f q_{x,f},\\
 q_{x,C}^*&=W_{x,C}^*/M_C^*,
 \end{aligned}
\]
where $\mathcal F_f$ is the time-integrated dry-mass flow, in kg, over the
interval used by this transport operation. The asterisk denotes the transported
endpoint, not the final source-adjusted state. A retained flux integrated over
one interval cannot be reused as though it represented a different interval
without an explicit conversion justified by the integrator.

Private moments need the same care as visible water. If a number is stored per
kg dry air, its inventory follows the dry-mass carrier. If it is stored per
volume, its numerical transport must use that definition. Moving mass with one
carrier while moving particle number with another can create a false change in
mean size even if both arrays remain positive. This affects terminal velocity,
collection and radiative interpretation at the next process call.

A basic invariant is the uniform-tracer test: a constant mixing ratio should
remain constant when only its carrier moves. Another is a closed or periodic
inventory test. A third is the boundary flux balance for an open domain. These
tests answer different questions and should not be replaced by a visual check
that a cloud retains roughly the same shape.

\section{A microphysics call is a coupled material transaction}

At a transported cell, microphysics receives pressure, temperature/theta,
dry density and the selected water and moment state. Processes redistribute
water among vapor, liquid and ice, change particle counts and sizes, exchange
latent energy, and move precipitation vertically. A process list is not a
set of independent increments that may all consume the original reservoir.
Ordering, shared-vapor allocation, collection routing and moment repair define
the actual closure.

Write the local phase-source inventory change as
\[
 \Delta W_x=M_C\Delta q_x.
\]
With no external water source and no sedimentation across the cell boundary,
the phase-change identity is $\sum_x\Delta W_x=0$. A domain water budget must
add sedimentation and other boundary exchanges. The appropriate total is all
prognostic water categories, not just cloud condensate or precipitation.
A cloud-to-rain conversion can reduce displayed cloud while conserving water.
An evaporating raindrop can disappear from the rain field while increasing
vapor and cooling the cell.

The thermodynamic increment follows the declared phase energies and heat
capacities. For an illustrative constant-pressure condensation of a small
amount $\delta q$ with constant latent heat and heat capacity,
\[
 \Delta T\simeq \frac{L_v\delta q}{c_p}.
\]
This is an explanatory approximation with held coefficients and the stated
thermodynamic constraint. The earlier microphysics chapter describes more
complete mixture and material-temperature issues. Freezing, melting, liquid
surface temperature, vapor competition and EOS projection prevent this
one-line estimate from being a universal energy owner.

A collection rule also makes a distributional assumption. Combining two mass
reservoirs without a compatible number change alters the inferred spectrum.
A count cap is therefore a model choice with consequences for size, fall speed
and subsequent collection, not merely protection against a large integer.
The documented Nimbus routing choices, including its retained new-hail gate
and count bounds, must be stated when comparing the model with a publication.
The guarded third-moment owners do not yet establish a complete live
source/EOS transaction; their source presence is not public readiness.

\section{Pressure adjustment is part of the thermal path}

The production dry rho--theta EOS can be expressed, with
$\kappa=R_d/c_p$ and $\gamma=1/(1-\kappa)$, as
\[
 p=p_{\rm ref}\left(\frac{\rho_dR_d\theta}{p_{\rm ref}}\right)^\gamma,
 \qquad T=\theta\left(\frac{p}{p_{\rm ref}}\right)^\kappa.
\]
These equations describe the declared dry pressure law. Moisture affects other
parts of the model, including phase energy and loading; that does not make the
pressure relation a full moist-mixture EOS.

The thermal sensitivity depends on what is held fixed. If pressure is held,
$dT/T=d\theta/\theta$. If dry density is held and pressure is subsequently
projected by the EOS, then
\[
 \frac{dT}{T}=\gamma\frac{d\theta}{\theta}.
\]
Thus a theta increment designed using a fixed-pressure temperature conversion
and a later fixed-density pressure projection form a different endpoint map
from the original fixed-pressure calculation. The equation is a local
sensitivity demonstration, not a claim that every Nimbus process applies that
simplified source or has an error of exactly $\gamma$.

To prescribe a particular endpoint temperature at held dry density, one can
instead solve the EOS relation for the required theta,
\[
 \theta=T^{1-\kappa}
        \left(\frac{p_{\rm ref}}{\rho_dR_d}\right)^\kappa.
\]
But choosing that endpoint requires an appropriate energy constraint first.
Replacing one conversion with this expression without deciding whether the
physical process is a constant-volume, constant-pressure or open-system
exchange would simply conceal a different assumption in the code.

Nimbus has explicit bounded work on a source/EOS energy mapping, described in
the preceding chapter. The live orchestration can place SGS, diffusion and
other operations between a microphysics source and final pressure projection.
Their changes must be represented in the complete endpoint transaction. A
cell-level mapping that rejects intervening mutation cannot be treated as a
complete field-level integration merely because its isolated residual is small.

\section{Mixing, sedimentation and radiation close different budgets}

Subgrid mixing exchanges momentum and scalar inventories between locations.
The eddy diffusivity depends on the closure and sometimes on gradients of the
state being mixed. A diffusivity has units of area per time; it is not itself
a mass flux. Density, face area, gradient and the integration interval turn it
into an exchange. A dissipative velocity update also needs a declared treatment
of its kinetic-energy change. An energetically complete closure cannot be
inferred solely from the sign of a diffusion coefficient.

Sedimentation is transport relative to the air. Its vertical flux depends on
the category's mass or moment distribution and its terminal-velocity closure.
The bottom boundary contributes precipitation leaving the atmosphere. If a
surface accumulation is updated without removing the corresponding atmospheric
inventory, water is created in the combined atmosphere-plus-surface accounting.
If removal occurs without the accumulation receipt, the atmospheric budget may
be correct while the published precipitation diagnostic is wrong.

For radiation, column heating follows flux divergence with the actual physical
layer mass and the chosen heat-capacity convention. The surface and top-of-column
radiative fluxes are external energy exchanges. Radiation should not be added
to a closed-domain energy expectation as though it were an internal transfer.
A held heating tendency and a freshly recomputed diagnostic flux can also refer
to different states; comparing them without their cadence information can
produce a misleading closure residual.

The combination of a boundary-layer scheme and resolved-grid SGS requires an
explicit division of transport. Both can represent unresolved mixing, but they
operate with different assumptions and length scales. Adding them together
without a coupling policy can count the same transport twice. Nimbus's
configuration admission and runtime owners restrict supported combinations;
a paper describing either component independently does not settle the combined
closure.

\section{A numerical ledger can be exact while the physics is wrong}

Let a completed step consist of states $X_0,X_1,\ldots,X_m$, and let
$E(X)$ be one consistently defined energy diagnostic. Then
\[
 E(X_m)-E(X_0)=\sum_{r=0}^{m-1}\big[E(X_{r+1})-E(X_r)\big]
\]
is a telescoping identity. It holds even if an operator implements an incorrect
physical law. Recording endpoint differences establishes where the diagnostic
changed, but does not independently establish why that change is physically
justified.

A stronger report compares each change with a separately specified exchange,
\[
 R_r=E(X_{r+1})-E(X_r)-\mathcal W_r-\mathcal Q_r-\mathcal B_r,
\]
where $\mathcal W_r$ is the declared mechanical work, $\mathcal Q_r$ the thermal
or phase exchange, and $\mathcal B_r$ the boundary contribution for that
operation and energy basis. Those terms must be derived from the appropriate
physics and geometry. Defining $\mathcal Q_r$ to equal the measured endpoint
difference by construction makes $R_r=0$ without testing the thermal law.

Floating-point storage is another observable operation. A double-precision
candidate cast back to a float field can change the diagnosed inventory.
Nimbus distinguishes storage/rounding effects where required by an owner;
those effects should not be relabeled as condensation, surface work or an
unknown physical sink. Deterministic reduction can make the effect repeatable,
but does not eliminate it.

The same discipline applies to conservation of water. A perfect domain sum can
hide a spatially wrong transfer between cells; a correct cell source can coexist
with a wrong boundary flux. Scientific validation therefore combines inventory
closure, local reference states, geometric invariants and dynamic behavior.
A single global residual is necessary for some claims and insufficient for all
of them.

\section{A worked control-volume example}

Consider an illustrative flat cell with volume
$100\,\mathrm m\times100\,\mathrm m\times50\,\mathrm m
=5\times10^5\,\mathrm{m^3}$ and dry density
$1\,\mathrm{kg\,m^{-3}}$. Its dry mass is $5\times10^5$ kg. At
$q_v=0.01\,\mathrm{kg\,kg^{-1}}$, it contains 5000 kg of vapor. These numbers
are chosen to expose the accounting; they are not a recommended grid or a
Nimbus regression reference.

Suppose one face brings in 1000 kg of dry air carrying the same vapor mixing
ratio while another removes 500 kg. The new dry mass is 500500 kg and vapor
mass is $5000+10-5=5005$ kg. Dividing gives the same $q_v=0.01$. Updating water
with the entering and leaving fluxes but dividing by the old dry mass would
instead produce a false concentration increase. Updating the mixing ratio
before constructing its old inventory can create the opposite kind of error.

Now let local condensation transfer $\delta q=2\times10^{-4}$ from vapor to
cloud liquid at the transported endpoint. It transfers 100.1 kg of water:
\[
 W_v'=5005-100.1,\qquad W_c'=W_c+100.1.
\]
Their sum is unchanged. With the illustrative held values
$L_v=2.5\times10^6\,\mathrm{J\,kg^{-1}}$ and
$c_p=1004\,\mathrm{J\,kg^{-1}\,K^{-1}}$, a constant-pressure estimate gives
$\Delta T\simeq0.498$ K. This estimate neglects temperature dependence,
mixture heat capacity, saturation feedback and all transport during the source.
Those are declared assumptions, not small corrections whose omission is
universally justified.

If part of that liquid converts to rain, the atmospheric total water still
need not change. If 20 kg subsequently falls through the physical bottom
boundary, atmospheric water decreases by 20 kg and the surface precipitation
inventory increases by 20 kg, subject to the chosen diagnostic units and area.
Rain falling from an upper cell into this cell is an internal vertical exchange,
not surface precipitation. The distinction depends on the physical interface
owner, especially on a stretched or terrain-following column.

Finally, pressure projection and momentum evolution determine the new thermal
and mechanical endpoint. The illustrative latent-heat estimate is not enough
to predict that endpoint without the EOS and work assumptions. A complete
step report would identify the carrier transport, phase transfer, sedimentation,
pressure-related energy change, mixing, boundaries and storage effects. The
example shows how to audit a step; it does not replace the actual source laws
or their tests.

\section{Publication should preserve the interpretation of the endpoint}

After the numerical state is accepted, diagnostics reconstruct quantities such
as temperature, vorticity, radar products and storm measures from that endpoint.
A diagnostic can use a different spatial location from its inputs: for example,
wind may be interpolated from faces to a scalar point for a plot. Its metadata
must state that reconstruction. A visual field must not silently combine the
new theta with an old pressure or a new geometry with an old domain origin.

The output layer publishes completion and geometry information so readers can
identify a consistent frame. Compression and visual downsampling are additional
representations of that frame. A block maximum is useful for retaining a small
visible peak but is not a conservative volume average. A volume renderer's
interpolated color is not a microphysics concentration. The architecture guide
explains these later representations and the GPU path that produces pixels.

Restart has a stronger requirement than visualization. It must preserve the
state needed by the next update, including private moments and relevant
histories, and must reject incompatible continuation assumptions. A successful
plot decode does not prove restart fidelity, and a successful restart does not
by itself prove that the underlying physical experiment is valid.

\section{Which assumptions travel with a result}

A reproducible scientific result should state the physical domain and frame,
scalar ownership and staggering, physical spacings and terrain geometry, base
state and initial perturbations, transport and time-integration choices,
boundary conditions, microphysics categories and moment units, closure and
radiation choices, source cadences and required coefficient identities. It
should also state precision, backend, threading and output transformations
when those affect the comparison. These are the components of the problem
actually solved.

The highest-level assumption is that a resolved, compressible grid model with
bulk closures is an appropriate description for the question being asked.
Subgrid distributions, unresolved turbulent motions and radiative optical
properties are modeled rather than resolved particle by particle. Refining the
grid does not remove every closure assumption. Conversely, matching one
observed storm feature does not uniquely validate each closure: compensating
errors can produce a similar visible outcome.

Verification asks whether the implemented operators satisfy their stated
mathematical contracts. Validation asks whether the model is adequate for a
physical use and whether its assumptions are supported by comparison with
independent evidence. Long campaigns examine behavior beyond a bounded case;
independent domain review examines the scientific interpretation. Nimbus's
readiness terminology keeps those stages separate. Engineering-incomplete
couplings remain work to be completed, even when a local formula, guard or
reference test is already correct.
